% Sumber: instrumen "Form Penggunaan dan Evaluasi Sistem Diabetes Clinical Decision
% Support" (Google Forms, 12 bagian), beserta berkas tanggapannya. Arsip instrumen dan
% berkas tanggapan ada pada docs/arsip/.
% Kode pasien pada kolom keempat adalah pasien yang WAJIB dipilih penilai pada tiap
% skenario, sehingga seluruh penilai menghadapi keadaan sistem yang sama.
\begin{table}[htbp]
\centering
\caption{Rancangan instrumen evaluasi ahli: dua belas bagian, sasaran yang diukur tiap bagian, dan kebutuhan yang hendak diuji olehnya. Kolom skala menyatakan rentang jawaban tertutup; bagian tanpa skala berisi isian bebas atau keterangan.}
\label{tab:instrumen-ahli}
\small
\begin{tabularx}{\textwidth}{@{}lLlcl@{}}
\toprule
\textbf{Bagian} & \textbf{Sasaran ukur} & \textbf{Pasien} & \textbf{Butir} & \textbf{Skala} \\
\midrule
1--3 & Persetujuan ikut serta, profil responden, dan petunjuk pengujian & -- & 6 & terbuka \\
\addlinespace
4 & Konsistensi prediksi dan status kondisi terhadap konteks pasien pada pemantauan CGM & P001 & 1 & 1--4 \\
\addlinespace
5 & Relevansi prediksi sekaligus kejelasan penyampaian keterbatasan data pada pemantauan \emph{finger-stick} & P007 & 1 & 1--4 \\
\addlinespace
6 & Mutu rekomendasi pada kondisi stabil, dinilai atas empat kriteria & P004 & 4 & 1--4 \\
\addlinespace
7 & Mutu rekomendasi pada kondisi hipoglikemia, beserta kesesuaiannya terhadap pedoman & P003 & 5 & 1--4 \\
\addlinespace
8 & Mutu rekomendasi pada kondisi hiperglikemia, beserta kesesuaiannya terhadap pedoman dan catatan bebas & P002 & 6 & 1--4 \\
\addlinespace
9 & Kejelasan penolakan memprediksi ketika riwayat tidak memadai, dan pengendalian kesan kepastian & P008 & 2 & 1--4 \\
\addlinespace
10 & Pengalaman penggunaan: sepuluh butir \emph{System Usability Scale} beserta lima butir mutu aplikasi & bebas & 15 & 1--5 \\
\addlinespace
11 & Penilaian menyeluruh: keterwakilan populasi sumber data, otonomi dokter dan bias, persetujuan dan keamanan data, kejelasan dasar rekomendasi, dan kelayakan alur pada praktik klinis di Indonesia & bebas & 5 & 1--4 \\
\addlinespace
12 & Kritik, saran, hambatan penerapan, dan risiko yang perlu diperhatikan & -- & 1 & terbuka \\
\bottomrule
\end{tabularx}

\vspace{0.4em}
{\tabnotestyle Empat kriteria pada bagian 6--8 adalah \emph{relevansi}, \emph{kejelasan}, \emph{kecukupan konteks}, dan \emph{batasan keputusan}; keempatnya dinilai pada satu kisi yang sama sehingga ketiga kondisi dapat dibandingkan berdampingan. Butir ``kesesuaian terhadap pedoman'' pada bagian 7 dan 8 bersifat kategoris dengan empat pilihan, yaitu sesuai pedoman dan aman diikuti, sesuai pedoman tetapi perlu verifikasi dokter, kurang sesuai pedoman dan berpotensi berisiko, serta tidak sesuai pedoman. Kode pasien pada kolom ketiga ditetapkan di muka dan wajib dipakai, sehingga seluruh penilai menghadapi keadaan sistem yang sama; tanpa pematokan itu, jawaban antar-penilai tidak dapat dibandingkan karena masing-masing akan menilai keluaran yang berbeda.\par}
\end{table}
